% title:          Free semigroup of symmetric integer matrices
% area:           matrices, free-groups, other-structures
% methods:
% difficulty:     hard
% difficulty-ai:  hard
% status:
% review:         human
% --- history: all optional, leave EMPTY if unknown ---
% origin:         journal
% origin-date:    2009
% origin-number:  Problem 11446
% also-in:
% taken-from:
% references:     American Mathematical Monthly, 2009, Problem 11446 (Vol. 116, No. 7, Aug.-Sep. 2009, per the legacy label), proposed by Christopher J. Hillar and Lionel Levine; the proposers' original statement - http://web.archive.org/web/20110915224316/http://www.msri.org/people/members/chillar/files/injectiveword.pdf ; listed as '[11446] Matrices whose products are all different' on Hillar's problems page - http://web.archive.org/web/20190824005320/http://www.msri.org/people/members/chillar/problems.html ; solution hints from the legacy file: ping-pong lemma - https://en.wikipedia.org/wiki/Ping-pong_lemma ; Mathlib Monoid.CoprodI.lift_injective_of_ping_pong - https://leanprover-community.github.io/mathlib4_docs/Mathlib/GroupTheory/CoprodI.html ; video 'Lemme du ping-pong et groupe modulaire' by Phil Caldero - https://youtu.be/2G41uyrIpws (matrices U and S); related classical fact that [[1,1],[0,1]] and [[1,0],[1,1]] generate a free monoid - https://math.stackexchange.com/questions/4177999
% history:        human
\begin{problem}
Prove or disprove: there exist \( 2 \times 2 \) symmetric integer matrices \( A \) and \( B \) such that no element of the multiplicative semigroup generated by \( A \) and \( B \) can be written in two different ways. (Thus, \( A, B, AA, AB, BA, BB, AAA, AAB, \ldots \) are all different.)
\end{problem}

%There is the jesus-christ author solution with the ping-pong lemma
%There is also the Spas solution with projective space (treat u,v as u/v, do equivalence class) and then try to have ABAA... be like a continued fraction, such that when you evaluate them, you get different numbers, like 1/(2+1/(1+...+x < 1/(2+1/(2+...+x for input vector x,1 (correxponding to the number x=x/1)
